
\documentclass[aps,prx,preprint,groupedaddress,showkeys]{revtex4-2}


\usepackage{amssymb}
\usepackage{amsmath}
\usepackage{graphicx}
\usepackage{relsize}
\usepackage{verbatim}
\usepackage[utf8]{inputenc} 
\usepackage{euscript}
\usepackage{amsfonts}


\newtheorem{defini}{Lemma }
\newtheorem{ansatz}{Ansatz}
\newcommand{\mds}[1] {\frac{d {#1}}{d\tau}}
\newcommand{\mdsd}[1] {\frac{d^2 {#1}}{d\tau^2}}
\newcommand{\mdt}[1] {\frac{d {#1}}{dt}}
\newcommand{\mdtd}[1] {\frac{d^2 {#1}}{dt^2}}
\newcommand{\mdf}[1] {\frac{d {#1}}{d\varphi}}


\begin{document}

\title{Developed Method. Interactions and their quantum picture.}
\author{Piotr Ogonowski}
\email{piotrogonowski@kozminski.edu.pl}
\affiliation{Kozminski University, Jagiellonska 57/59, Warsaw, 03-301, Poland}


\date{\today}

\begin{abstract}
By developing the previously proposed method of combining continuum mechanics with Einstein Field Equations, it has been shown that the classic relativistic description, curvilinear description, and quantum description of the physical system may be reconciled. For a system with an electromagnetic field, the Lagrangian density equal to the invariant of the electromagnetic field was obtained, vanishing four-divergence of canonical four-momentum appears to be consequence of the Poyinting theorem, and explicit form of one of gauges of the electromagnetic four-potential was introduced. The proposed method allows for further development with additional fields.
\end{abstract}

\keywords{Field theory, Lagrangian mechanics, Quantum mechanics, General Relativity, Unification of interactions}

\maketitle

\section{Introduction}

Over the past decades, great strides have been made in attempts to combine quantum description of interactions with General Relativity \cite{INT1}. There are currently many promising approaches to connecting the Quantum Mechanics and General Relativity, including perhaps the most promising ones: Loop Quantum Gravity \cite{LQG1}, \cite{LQG2}, \cite{LQG3}, String Theory \cite{ST1}, \cite{ST2}, \cite{ST3} and Noncommutative Spacetime Theory \cite{NS1}, \cite{NS2}. \\
~\\
A lot of work has also been done to clear up some challenges related to General Relativity and $\Lambda CMD$ model \cite{LCMD}. An explanation for the problem of dark energy \cite{DE1} and dark matter \cite{DM1} is still being sought, and efforts are still being made to explain the origin of the cosmological constant  \cite{CC1}, \cite{CC2}, \cite{CC3}.\\
~\\
The author also tries to bring his own contribution to the explanation of the above physics challenges, based on a recently discovered method, described in \cite{PMofC}. As this article will show, this method seems very promising and may help clarify at least some of the issues mentioned above. The author's method, similar to the approach presented in \cite{DFF}, \cite{LLM}, \cite{GEM1} also points to the essential connections between electromagnetism and Gereral Relativity, however, the postulated relationship is of a different nature and can be percieved as a generalization of the direction of research proposed in \cite{VGRA}, \cite{PGRA}, \cite{AGRA}.\\
~\\
According to conclusions from \cite{PMofC}, the description of motion in curved spacetime and its description in flat Minkowski spacetime with fields are equivalent, and the transformation between curved spacetime and Minkowski spacetime is known, because the geometry of curved spacetime depends on the field tensor. This allows for a significant simplification of research, because the results obtained in flat Minkowski spacetime can be easily transformed into curved spacetime. The last missing link seems to be the quantum description.\\
~\\
In this article, the author will focus on developing the method proposed in \cite{PMofC} in such a way, as to obtain the convergence of the description with the description of quantum mechanics. In the first chapter, the Lagrangian density for the system will be derived, allowing to obtain the tensor described in \cite{PMofC}. These conclusions will be used later in the article to propose quantum description of the system.\\
~\\ 
The author uses the Einstein summation convention, metric signature $(+,-,-,-)$ and commonly used notations. In order to facilitate the analysis of the article, the key conclusions from \cite{PMofC} are quoted in the subsection below. 

\subsection{Short summary of the method} \label{ShortSummary}

According to \cite{PMofC}, stress-energy tensor $T^{\alpha \beta}$ for a system with electromagnetic field in a given spacetime, described by a metric tensor $g^{\alpha \beta}$ is equal to
\begin{equation} \label{CH2}
T^{\alpha \beta} =  \varrho \, U^{\alpha} U^{\beta} - \left ( c^2 \varrho +  \Lambda_{\rho}  \right ) \left ( g^{\alpha \beta} - \xi \, h^{\alpha \beta} \right )
\end{equation}
where $\varrho_o$ is for rest mass density and
\begin{equation} \label{VRD}
\varrho \equiv \varrho_o \gamma
\end{equation}
\begin{equation} \label{CH4}
\frac{1}{\xi} \equiv \frac{1}{4}\, g_{\mu \nu} \, h^{\mu \nu}
\end{equation}
\begin{equation} \label{CH5}
\Lambda_{\rho} \equiv \frac{1}{4\mu_o} \mathbb{F}^{\alpha \mu} \, g_{\mu \gamma} \, \mathbb{F}^{ \beta \gamma } g_{\alpha \beta}
\end{equation}
\begin{equation} \label{CH6}
h^{\alpha \beta} \equiv 2 \; \frac{\mathbb{F}^{\alpha \delta} \, g_{\delta \gamma} \, \mathbb{F}^{ \beta \gamma } }{ \sqrt{\mathbb{F}^{\alpha \delta} \, g_{\delta \gamma} \, \mathbb{F}^{ \beta \gamma } \, g_{\mu \beta} \, \mathbb{F}_{\alpha \eta} \, g^{\eta \xi} \, \mathbb{F}^{\mu}_{\; \; \xi }}}
\end{equation}
where $\mathbb{F}^{\alpha \beta}$ represents electromagnetic field tensor, and where the stress–energy tensor for electromagnetic filed, denoted as $\Upsilon^{\alpha \beta}$ may be presented as follows
\begin{equation} \label{CH7}
\Upsilon^{\alpha \beta} \equiv  \Lambda_{\rho} \left (g^{\alpha \beta} - \xi \, h^{\alpha \beta}  \right ) = \Lambda_{\rho} g^{\alpha \beta} - \frac{1}{\mu_o}\mathbb{F}^{\alpha \delta} \, g_{\delta \gamma} \, \mathbb{F}^{ \beta \gamma }
\end{equation}
Thanks to the proposed amendment to the continuum mechanics, in flat Minkowski spacetime occurs
\begin{equation} \label{CH1}
\partial_{\alpha} U^{\alpha} = -\mdt{\gamma}  \quad \to \quad \partial_{\alpha} \, \varrho U^{\alpha} = 0
\end{equation}
thus denoting four-momentum density as $\varrho U^{\mu} = \varrho_o \gamma \, U^{\mu}$, total four-force density $f^{\mu}$ acting in the system is
\begin{equation} \label{TotFoI}
f^{\mu} \equiv \varrho A^{\mu} = \partial_{\alpha} \varrho U^{\mu} U^{\alpha}
\end{equation}
Denoting rest charge density in the system as $ \rho_o$ and
\begin{equation} \label{RHZERO}
\rho \equiv \rho_o \gamma
\end{equation}
electromagnetic four-current $J^{\alpha}$ is equal to
\begin{equation} \label{EN11}
J^{\alpha} \equiv \rho \, U^{\alpha} = \rho_o \gamma \, U^{\alpha}
\end{equation}
The pressure p in the system is equal to
\begin{equation} \label{Pressure}
p \equiv c^2 \varrho +  \Lambda_{\rho}
\end{equation}
In the flat Minkowski spacetime, total four-force density $f^{\alpha}$ acting in the system calculated from $\partial_{\beta} \; T^{\alpha \beta} =0$ is the sum of electromagnetic ($f^{\alpha}_{EM}$), gravitational ($f^{\alpha}_{gr}$) and the sum of remaining ($f^{\alpha}_{oth}$) four-force densities
\begin{equation} \label{forceDen}
f^{\alpha} = 
\begin{cases} 
f^{\alpha}_{EM} \equiv -\Lambda_{\rho} \, \partial_{\beta} \xi \, h^{\alpha \beta} \quad (electromagnetic) \\
+ \\
f^{\alpha}_{gr} \equiv \left ( g^{\alpha \beta} - \xi \, h^{\alpha \beta}  \right ) \partial_{\beta} \, p \quad (gravitational) \\
+ \\
f^{\alpha}_{oth} \equiv \frac{\varrho c^2}{\Lambda_{\rho}} f^{\alpha}_{EM}  \quad (sum \, of\, remaining \, forces)
\end{cases}
\end{equation}
As was shown in \cite{PMofC}, in curved spacetime ($g_{\alpha \beta} =h_{\alpha \beta}$) presented method reproduces Einstein Field Equations with an accuracy of $\frac{4\pi G}{c^4}$ constant and with cosmological constant $\Lambda$ dependent on invariant of electromagnetic field tensor $\mathbb{F}^{\alpha \gamma}$
\begin{equation} \label{CosmoCon}
\Lambda =  - \frac{\pi G}{c^4 \mu_o} \, \mathbb{F}^{\alpha \mu} \, h_{\mu \gamma} \, \mathbb{F}^{ \beta \gamma } h_{\alpha \beta} = - \frac{4\pi G}{c^4} \Lambda_{\rho}
\end{equation}
where $h_{\alpha \beta}$ appears to be metric tensor of the spacetime in which all motion occurs along geodesics and where $\Lambda_{\rho}$ describes vacuum energy density.\\
~\\
It was also shown, that in this solution, Einstein tensor describes the spacetime curvature related to vanishing in curved spacetime four-force densities $f^{\alpha}_{gr} + f^{\alpha}_{oth}$.\\
~\\
The proposed solution allows to add additional fields while maintaining its properties.

\section{Lagrangian density for the system} \label{SepOfFields}

Since for the considered method the transition to curved spacetime is known (based on electromagnetic field tensor), the rest of the article will focus on the calculations in the Minkowski spacetime with presence of electromagnetic field, where $\eta^{\alpha \beta}$ represents Minkowski metric tensor. \\
~\\
Using a simplified notation
\begin{equation} \label{simplfNot}
\mds{\, \ln{(p)}} = U_{\mu} \partial^{\mu} \ln{(p)} = U_{\mu} \frac{\partial^{\mu} p}{p}
\end{equation}
it can be seen that the four-force densities resulting from the obtained stress-energy tensor (\ref{forceDen}) in flat Minkowski spacetime can be written as follows
\begin{equation} \label{forceDenNew}
\begin{cases} 
f^{\alpha}_{gr} = \left ( \eta^{\alpha \beta} - \xi \, h^{\alpha \beta}  \right ) \partial_{\beta} p = \mds{\,\ln{(p)}} \varrho U^{\alpha} - T^{\alpha \beta}  \partial_{\beta} \ln{(p)}  \\
f^{\alpha}_{EM} = \frac{\Lambda_{\rho}}{p} \left ( f^{\alpha} - f^{\alpha}_{gr} \right )\\
f^{\alpha}_{oth} = \frac{\varrho c^2}{p} \left ( f^{\alpha} - f^{\alpha}_{gr} \right )
\end{cases}
\end{equation}
where $f^{\mu}_{EM}$ can also be represented in terms of electromagnetic four-potential and four-current. This means that to fully describe the system and derive the Lagrangian density, it is enough to find an explicit equation for the gravitational force or some gauge of electromagnetic four potential.\\
~\\
Referring to definitions from section \ref{ShortSummary} one may notice, that by proposing following electromagnetic four-potential $\mathbb{A}^{\mu}$
\begin{equation} \label{EM4pot}
\mathbb{A}^{\mu} \equiv - \frac{\Lambda_{\rho}}{p} \frac{\varrho_o}{\rho_o} U^{\mu}
\end{equation}
one obtains electromagnetic four-force density $f^{\alpha}_{EM}$ in form of
\begin{equation} \label{EM4EMf}
f^{\alpha}_{EM} = J_{\beta} \left ( \partial^{\alpha}\mathbb{A}^{\beta} - \partial^{\beta}\mathbb{A}^{\alpha} \right ) = \frac{\Lambda_{\rho}}{p} \left ( f^{\alpha} - \mds{\,\ln{(p)}} \varrho U^{\alpha} + \varrho c^2 \partial^{\alpha} \ln{(p)}  \right )
\end{equation}
where $J_{\beta}$ is electromagnetic four-current and where Minkowski metric property was utilized
\begin{equation} \label{MPR}
U_{\beta}U^{\beta} = c^2 \quad \rightarrow \quad U_{\beta}\partial^{\alpha}U^{\beta}=\frac{1}{2}\;  \partial^{\alpha} \left ( U_{\beta}U^{\beta} \right ) = 0 
\end{equation}
and where the forces in the system can be described by the following equality
\begin{equation} \label{Equality}
J_{\beta} \left ( \partial^{\alpha}\mathbb{A}^{\beta} - \partial^{\beta}\mathbb{A}^{\alpha} \right ) + \varrho U_{\beta} \left ( \partial^{\beta} \frac{\varrho c^2}{p} U^{\alpha} - \partial^{\alpha}\frac{\varrho c^2}{p} U^{\beta} \right ) = \varrho U_{\beta} \left ( \partial^{\beta} U^{\alpha} - \partial^{\alpha} U^{\beta} \right ) = f^{\alpha}
\end{equation}
Comparing (\ref{forceDenNew}) and (\ref{EM4EMf}) it is seen, that introduced electromagnetic four-potential yields 
\begin{equation} \label{For0tens}
0 = \left ( T^{\alpha \beta} - \varrho c^2 \; \eta^{\alpha \beta} \right ) \, \partial_{\beta} \ln{(p)}
\end{equation}
which is equivalent to imposing following condition on normalized stress-energy tensor
\begin{equation} \label{normCond}
0 = \partial_{\beta} \left ( \frac{T^{\alpha \beta}}{\eta_{\mu \gamma}T^{\mu \gamma}}  \right ) + \partial^{\alpha} \ln{\left ( \eta_{\mu \gamma} T^{\mu \gamma} \right )}
\end{equation}
and yields gravitational four-force density in Minkowski spacetime in form of
\begin{equation} \label{SIMGR}
f^{\mu}_{gr} = \varrho \left ( \mds{\, \ln{(p)}} U^{\mu} - c^2 \partial^{\mu} \ln{(p)} \right )
\end{equation}
Now, one may show, that proposed electromagnetic four-potential leads to correct solutions.\\
~\\
At first, recalling the classical Lagrangian density \cite{LagDen} for electromagnetism one may show why, in the light of the conclusions from \cite{PMofC} and above, it does not seem to be correct and thus makes it difficult to create a symmetric stress-energy tensor \cite{MEDP}. The classical value of the Lagrangian density for electromagnetic field, written with the notation used in the article, is
\begin{equation} \label{CLASSLAG}
- \mathcal{L}_{EMclassic} = \Lambda_{\rho} + \mathbb{A}^{\mu} J_{\mu}
\end{equation}
In addition to the obvious doubt that by taking the different gauge of the four-potential $\mathbb{A}^{\mu}$ one changes the value of the Lagrangian density, one may notice, that with considered electromagnetic four-potential, such Lagrangian density is equal to
\begin{equation} \label{CLASSLAG2}
- \mathcal{L}_{EMclassic} = \Lambda_{\rho} - \frac{\Lambda_{\rho}}{p} \frac{\varrho_o}{\rho_o} U^{\mu} \, U_{\mu} \rho_o \gamma = \Lambda_{\rho} - \frac{\Lambda_{\rho} \varrho c^2}{p} = \frac{\Lambda_{\rho}^2}{p}
\end{equation}
As it is seen, above Lagrangian density is not invariant under gradient over four-position and $\mathbb{A}^{\mu}$ and $J_{\mu}$ are dependent, what is not taken into account in classical calculation
\begin{equation} \label{DEPENDFIELDS}
\frac{\mathbb{A}^{\alpha}}{\mathbb{A}^{\mu} \mathbb{A}_{\mu}}  = \frac{J^{\alpha}}{\mathbb{A}^{\mu} J_{\mu}} 
\end{equation}
Above yields 
\begin{equation} \label{GRADPOT}
\frac{\partial \ln{\left ( \frac{1}{\sqrt{\mathbb{A}^{\mu} \mathbb{A}_{\mu}}} \right)}}{\partial \mathbb{A}_{\alpha}} = - \frac{J^{\alpha}}{\mathbb{A}^{\mu} J_{\mu}} = \frac{p}{\varrho c^2} \frac{J^{\alpha}}{\Lambda_{\rho}}
\end{equation}
One may decompose
\begin{equation} \label{DECPOT}
\ln{\left ( \frac{1}{\sqrt{\mathbb{A}^{\mu} \mathbb{A}_{\mu}}} \right)} = \ln{\left ( \frac{p \, \rho_o }{\Lambda_{\rho} \varrho_o c} \right)} = \ln{\left (p \right ) } - \ln{\left (\Lambda_{\rho} \right ) } - \ln{\left ( \frac{\varrho_o c}{\rho_o }\right ) } 
\end{equation}
and since $\frac{\varrho_o c}{\rho_o }$ are constants, one may simplify (\ref{GRADPOT}) to
\begin{equation} \label{GRAFIN}
\frac{\partial \ln{\left ( p \right)} }{\partial \mathbb{A}_{\alpha}} - \frac{\partial \ln{\left (\Lambda_{\rho} \right ) } }{\partial \mathbb{A}_{\alpha}} = \frac{p}{\varrho c^2} \frac{J^{\alpha}}{\Lambda_{\rho}} = \frac{J^{\alpha}}{\varrho c^2} + \frac{J^{\alpha}}{\Lambda_{\rho}}
\end{equation}
Above yields
\begin{equation} \label{MUINV}
\frac{\partial \ln{\left ( \Lambda_{\rho} \right ) }}{\partial \mathbb{A}_{\alpha}} = - \frac{J^{\alpha}}{\Lambda_{\rho}}
\end{equation}
which leads to the conclusion that $\Lambda_{\rho}$ acts as the Lagrangian density for the system
\begin{equation} \label{CLASSLAGNEW}
\frac{\partial \Lambda_{\rho}} {\partial \mathbb{A}_{\alpha}} =  \partial_{\nu}  \left (  \frac{\partial \Lambda_{\rho}}{\partial (\partial_{\nu} \mathbb{A}_{\alpha} )}   \right ) = - J^{\alpha}
\end{equation}
which would support conclusions from \cite{ALL1}, and thus stress-energy tensor for the system is equal to
\begin{equation} \label{TensSep1}
T^{\alpha \beta} = \frac{1}{\mu_o}\mathbb{F}^{\alpha \gamma} \partial^{\beta} \mathbb{A}_{\gamma} -  \Lambda_{\rho} \eta^{\alpha \beta}
\end{equation}
In fact, the proof of correctness of the electromagnetic field tensor (noted as $\Upsilon^{\alpha \beta}$) allows to see this solution
\begin{equation} \label{Ups1}
f_{EM}^{\beta} = \partial_{\alpha} \Upsilon^{\alpha \beta} = J^{\gamma}\mathbb{F}^{\beta}_{\, \, \, \gamma} - \frac{1}{\mu_o}\mathbb{F}^{\alpha \gamma} \partial_{\alpha} \mathbb{F}^{\beta}_{\, \, \, \gamma}
\end{equation}
what yields following property of electromagnetic field tensor
\begin{equation} \label{Ups2}
\mathbb{F}^{\alpha \gamma} \partial_{\alpha} \partial_{\gamma} \mathbb{A}^{\beta} 
 = \mathbb{F}^{\alpha \gamma}  \partial^{\beta} \partial_{\alpha} \mathbb{A}_{\gamma} \end{equation}
Using the above substitution, one may note
\begin{equation} \label{TensorProof}
\partial_{\alpha} \Upsilon^{\alpha \beta} = \partial_{\alpha} \frac{1}{\mu_o} \mathbb{F}^{\alpha \gamma} \partial_{\gamma} \mathbb{A}^{\beta} - \partial_{\alpha} \frac{1}{\mu_o} \mathbb{F}^{\alpha \gamma} \partial^{\beta} \mathbb{A}_{\gamma} = \frac{1}{\mu_o} \mathbb{F}^{\alpha \gamma} \partial^{\beta} \partial_{\alpha} \mathbb{A}_{\gamma} - J^{\gamma}\partial_{\gamma} \mathbb{A}^{\beta} - \partial_{\alpha} \frac{1}{\mu_o} \mathbb{F}^{\alpha \gamma} \partial^{\beta} \mathbb{A}_{\gamma}
\end{equation}
Therefore the invariance of $\Lambda_{\rho}$ with respect to $\mathbb{A}_{\alpha}$ and  $\partial_{\nu} \mathbb{A}_{\alpha}$ is both a condition on the correctness of the electromagnetic stress-energy tensor and on $\Lambda_{\rho}$ in the role of Lagrangian density
\begin{equation} \label{TENSCORR}
0 = \frac{\partial \Lambda_{\rho}}{\partial (\partial_{\alpha} \mathbb{A}_{\gamma})} \partial^{\beta} (\partial_{\alpha} \mathbb{A}_{\gamma}) + \frac{\partial \Lambda_{\rho}} {\partial \mathbb{A}_{\gamma}} \partial^{\beta} \mathbb{A}_{\gamma} = \frac{1}{\mu_o} \mathbb{F}^{\alpha \gamma} \partial^{\beta} \partial_{\alpha} \mathbb{A}_{\gamma}  - J^{\gamma} \partial^{\beta} \mathbb{A}_{\gamma} = \partial_{\alpha} \frac{1}{\mu_o} \mathbb{F}^{\alpha \gamma} \partial^{\beta} \mathbb{A}_{\gamma}
\end{equation}
what yields for (\ref{TensorProof}) that
\begin{equation} \label{UpsFin}
\partial_{\alpha} \Upsilon^{\alpha \beta} = J^{\gamma} \partial^{\beta} \mathbb{A}_{\gamma} - J^{\gamma}\partial_{\gamma} \mathbb{A}^{\beta} = f_{EM}^{\beta}
\end{equation}
Equations (\ref{CH2}), (\ref{CH7}) and (\ref{TensSep1}) yield
\begin{equation} \label{TensSep2}
\frac{1}{\mu_o} \mathbb{F}^{\alpha \gamma} \partial_{\gamma} \mathbb{A}^{\beta} = \varrho \, U^{\alpha} U^{\beta} - \frac{c^2 \varrho }{\Lambda_{\rho}} \, \Upsilon^{\alpha \beta}
\end{equation}
what yields second representation of the stress-energy tensor
\begin{equation} \label{TensSep3}
T^{\alpha \beta} = \frac{p}{\varrho c^2} \cdot \frac{1}{\mu_o} \mathbb{F}^{\alpha \gamma} \partial_{\gamma} \mathbb{A}^{\beta} - \frac {\Lambda_{\rho}}{c^2} U^{\alpha} U^{\beta} = \frac{p}{\varrho c^2} \, \partial_{\gamma} \frac{1}{\mu_o} \mathbb{F}^{\alpha \gamma} \mathbb{A}^{\beta}
\end{equation}
After four-divergence, it gives additional expression for relation between forces and gives useful clues about the behavior of the system when transitioning to the description in curved spacetime.

\section{Hamiltonian density and quantum picture} \label{Quantas}

Since 
\begin{equation} \label{HQ1}
\frac{\partial \Lambda_{\rho}}{\partial (\partial \mathbb{A}_{\gamma} / \partial x_{\alpha})}  = \frac{1}{\mu_o} \mathbb{F}^{\alpha \gamma} \quad \to \quad \frac{\partial \Lambda_{\rho}}{\partial (\partial \mathbb{A}_{\gamma} / \partial x_0)} = \frac{1}{\mu_o} \mathbb{F}^{0 \gamma}
\end{equation}
thus noting Hamiltonian density in terms of the Lagrangian density \cite{LDEM} (and also from (\ref{TensSep1})) one gets
\begin{equation} \label{HQ2}
\mathcal{H} = \left ( \frac{\partial \Lambda_{\rho}}{\partial (\partial \mathbb{A}_{\gamma} / \partial x_0)} \right ) \frac{\partial \mathbb{A}_{\gamma}}{\partial x_0} - \Lambda_{\rho} = \frac{1}{\mu_o c} \mathbb{F}^{0 \gamma} \, \frac{\partial \mathbb{A}_{\gamma}}{\partial t} - \Lambda_{\rho}
\end{equation}
where conjugate momentum field $\Pi^{\gamma}$ is
\begin{equation} \label{HQ3}
\Pi^{\gamma} \equiv \frac{1}{\mu_o} \mathbb{F}^{0 \gamma}
\end{equation}
As it is seen, above Hamiltonian density agrees with the classical Hamiltonian density for electromagnetic field \cite{LDEM} except that this Hamiltonian density was currently mainly considered for sourceless regions. According to the result above, this Hamiltonian density describes also gravitational interaction and other forces resulting from electromagnetism. Above, therefore, significantly simplifies quantum field theory equations  \cite{QFT1}, \cite{QFT2}, \cite{QFT3}.\\
~\\
At first one may notice, that in transformed (\ref{TensSep1})
\begin{equation} \label{HQ4}
- T^{\alpha 0} =  - \frac{1}{\mu_o} \mathbb{F}^{\alpha \gamma} \partial_{\gamma} \mathbb{A}^{0} + \Upsilon^{\alpha 0}
\end{equation}
first row of electromagnetic stress-energy tensor $\Upsilon^{\alpha 0}$ is a four-vector, representing energy density of electromagnetic field and Poynting vector \cite{PoV} - the Poynting four-vector. Therefore vanishing four-divergence of the $T^{\alpha 0}$ must represent Poynting theorem. Indeed, properties (\ref{Ups2}) and (\ref{TENSCORR}) provide such equality
\begin{equation} \label{HQ5}
0 = - \partial_{\alpha} T^{\alpha 0} = J_{\gamma} \mathbb{F}^{0 \gamma} + \partial_{\alpha} \Upsilon^{\alpha 0} 
\end{equation}
which also shows, that for the constant total energy of the system, $T^{\alpha 0}$ is independent of the four-position.\\
~\\
Next, one may introduce auxiliary variable $\varepsilon$ with the energy density dimension, defined as follows
\begin{equation} \label{EpsDef}
c\varepsilon \equiv - \frac{1}{\mu_o} \mathbb{F}^{0 \mu} \mds {\mathbb{A}_{\mu}}
\end{equation}
and comparing the result
\begin{equation} \label{Lagr}
- U_{\beta}T^{0 \beta} = c\varepsilon + c\gamma \Lambda_{\rho}
\end{equation}
between the two tensor definitions (\ref{TensSep1}), (\ref{TensSep3}) one may notice, that it must hold
\begin{equation} \label{EpsProp}
- \frac{p}{\varrho c^2} \cdot \frac{1}{\mu_o} \mathbb{F}^{0 \mu} \partial_{\mu} \mathbb{A}^{\beta} = - \frac{p}{\varrho c^2 \mu_o} \cdot \left ( U^{\beta} \, \mathbb{F}^{0 \mu} \partial_{\mu} \frac{\mathbb{A}^0}{c\gamma} + \frac{\mathbb{A}^0}{c\gamma} \mathbb{F}^{0 \mu} \partial_{\mu} U^{\beta} \right )  = \frac{\varepsilon}{c} \, U^{\beta} - \frac{p}{\varrho c^2} \frac{\mathbb{A}^0}{\gamma c\mu_o} \mathbb{F}^{0 \mu} \partial_{\mu} U^{\beta}
\end{equation}
because the second component of above vanishes contracted with $U_{\beta}$, due to the property of the Minkowski metric (\ref{MPR}). Therefore (\ref{TensSep1}) also yields
\begin{equation} \label{TensSepTrans}
- T^{0 \beta} =  \varepsilon \frac{\varrho c}{p} \, U^{\beta} - \frac{c \epsilon_o \mathbb{A}^0}{\gamma} \mathbb{F}^{0 \mu} \partial_{\mu} U^{\beta} + \Upsilon^{0 \beta}
\end{equation}
where $\epsilon_o$ is electric vacuum permittivity.\\
~\\
It now may be concluded, that after integration of the $\frac{1}{c}T^{0 \beta}$ with respect to the volume, the total energy transported in the isolated system should be the sum of the four-momentum and four-vectors describing energy transport related to fields. This would be consistent with the conclusion from \cite{SETH} that "equations of motion for matter do not need to be introduced separately, but follow from the field equations". In considered approach that would mean, that the cannonical four-momentum density is a part of the stress-energy tensor.\\
~\\
Therefore, by analogy with the Poyting four-vector $\frac{1}{c}\Upsilon^{0 \beta}$, one may introduce a four-vector $Z^{\beta}$ understood as its equivalent for the remaining interactions and rewrite (\ref{TensSepTrans}) as
\begin{equation} \label{HQ4b}
- \frac{1}{c} T^{0 \beta} =  \varrho_o U^{\beta} + Z^{\beta} + \frac{1}{c}\Upsilon^{0 \beta}
\end{equation}
where
\begin{equation} \label{HQ4d}
Z^{\beta} \equiv \left (\frac{\varepsilon \gamma}{p} - 1 \right )  \varrho_o U^{\beta} - \frac{\epsilon_o \mathbb{A}^0}{\gamma} \mathbb{F}^{0 \mu} \partial_{\mu} U^{\beta}
\end{equation}
~\\
Above also drives to conclusion, that
\begin{equation} \label{LagDenH}
\frac{1}{c\gamma} U_{\beta} T^{0 \beta} = - \frac{\varepsilon}{\gamma} - \Lambda_{\rho}
\end{equation} 
may act as a Lagrangian density, used in the classic relativistic description based on four-vectors.\\
~\\
Based on the knowledge of the existing Lagrangians, a probable assumption can be made at this point that
\begin{equation} \label{HQ4e}
\varepsilon = \varrho_o c^2
\end{equation}
what would yield
\begin{equation} \label{LDR}
- \frac{1}{c} U_{\beta} T^{0 \beta} = \varrho_o c^2 + \gamma \Lambda_{\rho}
\end{equation} 
\begin{equation} \label{HQ4f}
Z^{\beta} = \rho_o \mathbb{A}^{\beta} - \frac{\epsilon_o \mathbb{A}^0}{\gamma} \mathbb{F}^{0 \mu} \partial_{\mu} U^{\beta}
\end{equation}
where $- \frac{\epsilon_o \mathbb{A}^0}{\gamma} \mathbb{F}^{0 \mu} \partial_{\mu} U^{\beta}$, due to its properties, may be associated with some description of the spin.\\
~\\
The above result should not be surprising, because considering the system with electromagnetic field and cannonical four-momentum density as part of the stress-energy tensor, actually one of components of cannonical four-momentum density should be related to electromagnetic four-potential and charge density.\\
~\\
Finally, one may define another gauge $\bar{\mathbb{A}}_{\gamma}$ of electromagnetic four-potential $\mathbb{A}_{\gamma}$ in following way
\begin{equation} \label{HQ6}
\bar{\mathbb{A}}^{\gamma} \equiv  \mathbb{A}^{\gamma} - \partial^{\gamma} \mathbb{A}^{\beta}X_{\beta} = - X_{\beta} \partial^{\gamma} \mathbb{A}^{\beta}
\end{equation}
and note, that
\begin{equation} \label{HQ7}
- X_{\beta} T^{0 \beta} =  \frac{1}{\mu_o} \mathbb{F}^{0 \gamma} \bar{\mathbb{A}}_{\gamma} + X_{\beta} \Upsilon^{0 \beta}
\end{equation}
Since for constant total energy of the system
\begin{equation} \label{HQ8}
\partial^{\alpha} X_{\beta} T^{0 \beta} =  T^{0 \alpha}
\end{equation}
this brings two more important insights: 
\begin{itemize}
  \item $\frac{1}{c}X_{\beta} T^{0 \beta}$ may play a role of the density of Hamilton's principal function,
  \item Hamilton's principal function may be expressed based on the electromagnetic field only, so in the absence of the electromagnetic field it disappears.
\end{itemize}
~\\
One may now summarize the above and propose a method for quantizing the system and for the description of the system with the use of classical field theory.\\
~\\
At first, it should be noted, that the above reasoning changes the interpretation of what the relativistic principle of least action means. As one may conclude from above, there is no inertial system in which no fields act, and in the absence of fields, the Lagrangian, the Hamiltonian and Hamilton's principal function vanish. Since the metric tensor (\ref{CH6}) for description in curved spacetime depends on the electromagnetic field tensor only, it seems clear, that in the considered system, the absence of the electromagnetic field means actually the disappearance of spacetime and the absence of any action.\\
~\\
One may then introduce generalized, canonical four-momentum $H^{\mu}$ as four-gradient on Hamilton's principal function S
\begin{equation} \label{Q2}
H^{\mu} \equiv - \frac{1}{c} \int T^{0 \mu} \, d^3 x \equiv -\partial^{\mu} S
\end{equation}
where 
\begin{equation} \label{Q2a}
- S \equiv H^{\mu}X_{\mu}
\end{equation}
and where $H^{\mu}$ does not depend on four-position. One may also conclude from previous findings, that canonical four-momentum should be in form of
\begin{equation} \label{Q5}
H^{\mu} = P^{\mu} + V^{\mu} 
\end{equation}
where 
\begin{equation} \label{CD1}
V^{\mu} \equiv \int Z^{\mu} + \frac{1}{c} \Upsilon^{0 \mu} \, d^3 x
\end{equation}
and where four-momentum $P^{\mu}$ may be now considered as just "other gauge" of  $-V^{\mu}$ 
\begin{equation} \label{Q5b}
\partial^{\alpha} P^{\mu} = - \partial^{\alpha} V^{\mu}
\end{equation}
Since in the limit of the inertial system one gets $P^{\mu} X_{\mu} = mc^2\tau$, therefore, to ensure vanishing Hamilton's principal function in the inertial system, one may expect that
\begin{equation} \label{Q1}
V^{\mu} X_{\mu} \equiv - mc^2 \tau 
\end{equation}
what also yields vanishing in the inertial system Lagrangian $L$ in form of
\begin{equation} \label{Q4}
- \gamma L = U_{\mu} H^{\mu}= F^{\mu} X_{\mu} 
\end{equation}
where $F^{\mu}$ is four-force, and where all above yields
\begin{equation} \label{Q6}
V^{\mu} U_{\mu} - F^{\mu}X_{\mu} = - mc^2
\end{equation}
thus
\begin{equation} \label{Q6b}
H^{\mu} H_{\mu} + m^2c^2 = 2m F^{\mu}X_{\mu} + V^{\mu} V_{\mu}
\end{equation}
~\\
In this picture, the Hamilton's principal function, generalized canonical four-momentum and Lagrangian vanish for inertial system as expected. Assuming that equation (\ref{LDR}) is indeed correct and denoting
\begin{equation} \label{Elambda}
E_{\Lambda} \equiv - \int \Lambda_{\rho} \, d^3 x
\end{equation}
one obtains from (\ref{Q6b})
\begin{equation} \label{Q6c}
H^{\mu} H_{\mu} - m^2c^2 = V^{\mu} V_{\mu} - 2m\gamma E_{\Lambda}
\end{equation}
~\\
To ensure compatibility with the equations of quantum mechanics it suffices to assume that right side of above equation vanishes
\begin{equation} \label{Q7}
\gamma E_{\Lambda} = \frac{V^{\mu} V_{\mu}}{2m} 
\end{equation}
~\\
Thanks to above, by introducing quantum wave function $\Psi$ in form of 
\begin{equation} \label{Q8}
\Psi \equiv e^{\pm iK^{\mu} X_{\mu}}
\end{equation} 
where $K^{\mu}$ is wave four-vector related to cannonical four-momentum 
\begin{equation} \label{Q8b}
\hbar  K^{\mu} \equiv H^{\mu}
\end{equation}
from (\ref{Q6c}) one obtains Klein-Gordon equation
\begin{equation} \label{Q9}
\left ( \Box + \frac{m^2c^2}{\hbar^2} \right ) \Psi = 0
\end{equation} 
which allows for further analysis of the system in the quantum approach, eliminating the problem of negative energy appearing in solutions \cite{KGE}. \\
~\\
The above representation allows the analysis of the system in the quantum approach, classical approach based on (\ref{HQ4}) and the introduction of a field-dependent metric in (\ref{CH6}) for curved spacetime, which connects previously divergent descriptions of physical systems.

\section{Conclusions and Discussion} \label{3Chap}

As shown above, the proposed method summarized in section \ref{ShortSummary} seems to be very promising area of farther research. In addition to the earlier agreement with Einstein Field Equations in curved spacetime, by imposing condition (\ref{For0tens}) on normalized stress-energy tensor in flat Minkowski spacetime, one obtains consistent results, developing the knowledge of the physical system with electromagnetic field:
\begin{itemize}
  \item Lagrangian density for the systems appears to be equal to $\mathcal{L} = \Lambda_{\rho} = \frac{1}{4\mu_o} \mathbb{F}^{\alpha \beta} \mathbb{F}_{\alpha \beta} $
  \item Stress-energy tensor may be simplified to familiar form: $T^{\alpha \beta} = \frac{1}{\mu_o}\mathbb{F}^{\alpha \gamma} \partial^{\beta} \mathbb{A}_{\gamma} -  \Lambda_{\rho} \eta^{\alpha \beta}$
  \item $H^{\beta} \equiv - \frac{1}{c} \int T^{0 \beta} \, d^3 x$ acts as cannonical four-momentum for the point-like particle
  \item The vanishing four-divergence of $H^{\beta}$ turns out to be the consequence of Poynting theorem
  \item Some gauge of electromagnetic four-potential may be expressed as $\mathbb{A}^{\mu} = - \frac{\Lambda_{\rho}}{p} \frac{\varrho_o}{\rho_o} U^{\mu}$
  \item Gravitational, electromagnetic and sum of other forces acting in the system may be expressed as shown in (\ref{forceDenNew}) where gravitational four-force is dependent on the pressure $p$ in the system as shown in (\ref{SIMGR}) 
\end{itemize}
~\\
Obtained canonical four-momentum $H^{\mu}$ is equal to  
\begin{equation} \label{Q2}
H^{\mu} = P^{\mu} + V^{\mu} 
\end{equation}
where $P^{\mu}$ is four-momentum and $V^{\mu}$ describes the transport of energy due to the field. It is calculated as
\begin{equation} \label{QV1}
V^{\mu} = q \mathbb{A}^{\mu} + \mathbb{S}^{\mu} + Y^{\mu} 
\end{equation}
where $\mathbb{A}^{\mu}$ is electromagnetic four-potential, $\mathbb{S}^{\mu}$ due to its properties, seems to be some description of the spin, and where $Y^{\mu}$ is the volume integral of the Poyinting four-vector. \\
~\\
It also seems, that this approach allows to combine previously divergent methods of curvilinear and classic description with the quantum description. The proposed method should facilitate further research on the quantum picture of individual fields, significantly simplifying equations of the Quantum Field Theory and leading in natural way to Klein-Gordon equation (\ref{Q9}).\\
~\\
Further analysis using the variational method may also lead to next discoveries regarding both the theoretical description of quantum fields and elementary particles associated with them, and the possibility of experimental verification of the obtained results.

\section{Statements} \label{4Chap}

Data sharing is not applicable to this article, as no datasets were generated or analyzed during the current study.\\
~\\
The author did not receive support from any organization for the submitted work.\\
~\\
The author has no relevant financial or non-financial interests to disclose.\\


\section*{References}

\bibliography{DevMethodAPS}

\end{document}



